\documentclass[12pt]{article}

\usepackage[utf8]{inputenc}
\usepackage[T1]{fontenc}

\title{La lingüística de la biología: bases cognitivas del lenguaje y construcción del discurso científico}
\author{J. Toledo, S. Castro, A. Jara, D. Rojas}
\date{\today}

\begin{document}

\maketitle

\section{Introducción}
El lenguaje es una capacidad biológica de los seres humanos y a la vez, la lingüística sirve para estudiar como se escribe y se comunica la biologia. Así entendemos mejor tanto nuestra mente al comunicarnos como la forma en que se construye el saber biologico.

Las interacciones humanas han sido fundamentales para la organización y el desarrollo de las sociedades. Sin embargo esta cooperacion no depende unicamente de la proximidad entre individuos sino tambien de la capacidad para coordinar acciones, transmitir informacion y compartir conocimiento. Por ello aunque el ser humano no se distingue por su fuerza fisica, su capacidad cognitiva y su habilidad para establecer sistemas complejos de comunicacion es lo que le ha permitido transformar el entorno y desarrollar una organización social más abanzada.

Durante este proceso las formas de comunicación linguistica se volvieron pogresivamente más complejas hasta desarrollar sistemas lingüísticos más complejos. Por ello, el lengaje no se debe entender únicamente como un producto culturar destinado a facilitar la interaccion social, sino tambien como una capacidad biologica y cognitiva relacionada con la organizacion del cerebro, el procesamiento de informacion, y la posbilidad de ver el mundo de maneras diferentes. Esta perspectiva permite analizar el lenguaje como objeto de estudio de la biologia especialmente en relacion con su origen, su desarrollo y las restricciones cognitivas que intervienen en su funcionamiento

La biolingúística busca analizar la emergencia del lenguaje, no solo a partir de la presión comunicativa, sino también por principios de eficiencia y desarrollo cognitivo. La adquisición y el desarrollo del lenguaje pueden explicarse mediante la interacción entre la dotación biológica del individuo, la experiencia y ciertos principios cognitivos generales. Estos últimos constituyen el llamado tercer factortercer factor. En la propuesta generativista, este termino se refiere a principios que no son exclusivos del lenguaje, pero que pueden influir en su organización, como la reducción del esfuerzo necesario para procesar información. Por tanto, el lenguaje se desarrolla a partir de la interacción entre las capacidades biológicas del organismo, la experiencia y principios cognitivos más generales(Alconchel M. 2004).

.Los principios de eficiencia permiten comprender la relación entre la organización interna del lenguaje y su uso comunicativo. El sistema lingüístico debe equilibrar la economía del procesamiento interno con la necesidad de producir expresiones interpretables dentro de una situación comunicativa. Pero, para lograr el avance colectivo, es necesario compartir la información que se procesa internamente, por lo que debe expresarse de tal manera que pueda ser interpretada sin que pierda sentido. En consecuencia, existe una relación entre dos exigencias: reducir el esfuerzo cognitivo y hacer que el mensaje resulte comprensible en una situación comunicativa concreta. El lenguaje no responde de manera absoluta a una sola de estas exigencias, sino que busca establecer un equilibrio entre la eficiencia del procesamiento y la eficacia de la comunicación.

Esta relación puede analizarse mediante diferentes niveles de la lingüística. La microsintaxis estudia la estructura interna de la oración, incluyendo las relaciones entre sus componentes, la concordancia, las dependencias y el papel del verbo. La macrosintaxis, en cambio, analiza cómo se organizan unidades superiores a la oración para construir un mensaje completo, teniendo en cuenta la relación entre los enunciados, el contexto y la intención comunicativa. La pragmática estudia cómo se interpreta un mensaje de acuerdo con la situación en que se produce, los conocimientos compartidos entre los participantes y la finalidad de la comunicación.

La importancia de esta relación se observa especialmente en el lenguaje científico de la biología. La biología no solo produce conocimientos sobre los seres vivos, sino que también necesita expresarlos mediante conceptos, definiciones, descripciones, explicaciones y argumentos. Por esta razón, el análisis lingüístico permite estudiar cómo se organizan los textos biológicos y cómo se garantiza que la información científica sea clara, coherente e interpretable. En este contexto, la macrosintaxis y la pragmática contribuyen a comprender la organización global del mensaje, mientras que la microsintaxis permite examinar la estructura de las oraciones y las relaciones gramaticales que sostienen el discurso.

Referencias
Rodríguez, C. & Lara, E. (2017). Macrosintaxis y lingüística pragmática. Círculo de lingüística aplicada a la comunicación, 71, 5-34.
Alconchel, M. Á. (2004). Chomsky la gramática generativa. *Revista Digital "Investigación y Educación"*, 3(7). ISSN 1696-7208
Vitral, L.(2017). La emergencia de la biolingüística y sus consecuencias para la teoría lingúística. Lingüística y literatura.
ARCHIVO SIN ACTIVIDAD
\end{document}
